\documentclass[11pt]{article}

\usepackage[margin=1in]{geometry}
\usepackage{amsmath,amssymb}
\usepackage{booktabs,tabularx}
\usepackage[colorlinks=true,allcolors=blue]{hyperref}

\title{Response to Reviewer\\
\large \emph{Toward an Algorithmic Free Energy Principle for Explanatory Programs}}
\author{}
\date{\today}

\begin{document}

\maketitle

We thank the reviewer for a careful and constructive report. The revised manuscript addresses the concerns raised, with the exception of empirical validation (reviewer Major Concern~2), which will be reported in a companion paper. We have applied the changes in a phased manner following an internal revision plan; here we walk through each reviewer item, describe the change made, and point to the relevant section or paragraph of the revised manuscript.

\section*{Summary of Changes}

The revised manuscript differs from the original in the following principal ways.

\begin{itemize}
\item The Related Work section has been reorganized into four clusters, each closed by an explicit delta statement identifying what this paper contributes beyond the cluster it is closing. Substantial new engagement with information-theoretic bounded rationality (Ortega \& Braun), Bayesian experimental design (Lindley; Chaloner \& Verdinelli; MacKay), compression-based curiosity (Schmidhuber), universal-prediction convergence (Leike \& Hutter), information bottleneck (Bialek-Nemenman-Tishby; Tishby-Pereira-Bialek), and program induction (DreamCoder, ARC-AGI, FunSearch) has been added. A new Positioning subsection (Section~3.5) states the paper's stance relative to these literatures.
\item Deutsch's notion of hard-to-vary explanations is now explicitly cited at its first appearance (Section~2, criterion~4).
\item A new Section~8 discusses reference-machine invariance, posterior consistency, and exploration sufficiency. The ``normatively complete'' language has been replaced throughout by ``variationally specified at every finite history.''
\item A new Section~10 (The Open Mechanistic Challenge) replaces the previous brief sketch with a concrete compliance protocol specifying three internal quantities to be exposed, the verification test applied to each, the statistical-power analysis, and the interpretability substrates that could support the exposures.
\item Eight minor technical items have been tightened: uniqueness justifications, admissibility and zero-support handling, a Semimeasure-Completion Lemma (Section~7.1), explicit MAP existence, semantic-MDL well-definedness, motivation for the preference distribution $\rho_t$, effective enumerability of $\mathcal{P}$, and disambiguated notation for the joint/marginal predictive laws ($\Qcal_t^{PO}$ vs.\ $\Qcal_t^O$).
\item The Abstract, Introduction's ``sharp move'' paragraph, and Conclusion have been rewritten to temper novelty claims, consistent with the expanded Related Work.
\item Bibliography hygiene: the header-comment TODO block and the in-text ``NOTE TO AUTHOR'' comment have been removed. The bib file now contains all referenced entries.
\end{itemize}

\section*{Response to Major Concerns}

\paragraph{MC1 (novelty overstated; literature omissions).} Addressed. Section~3 now treats four clusters with explicit deltas. Ortega \& Braun, Lindley, Chaloner \& Verdinelli, MacKay, Schmidhuber, Bialek-Nemenman-Tishby, Still-Crutchfield, Leike-Hutter, DreamCoder (Ellis et al.), and recent program-induction benchmarks are all discussed. The Positioning subsection (3.5) is new and states the paper's position compactly. The Abstract and the ``sharp move'' paragraph of the Introduction have been rewritten to avoid the ``missing synthesis'' framing; we now describe the contribution as ``one particular assembly of components that have been circulating separately, with specific analytic and methodological payoffs.''

\paragraph{MC2 (no empirical content).} Partially addressed, partially deferred. We do not include empirical results in this revision; we agree that empirical validation of the behavioral proxy is essential, and we are preparing a companion paper that reports results on the piecewise-modular running example and a broader DSL-based rule-recovery benchmark. The present revision fixes all of the proxy parameters (Table~2, the complexity surrogate, the query penalty, and the held-out evaluation protocol) so that the empirical paper can report results against a stable target. Section~9.1 discusses this explicitly.

\paragraph{MC3 (``principled prior'' claim too strong).} Addressed. Section~8.1 states the invariance theorem precisely, acknowledges that $c_{U,V}$ can be large on any fixed finite-sample problem, and explicitly tempers the ``principled vs.\ stipulated'' framing: ``Replacing a stipulated FEP prior with the Solomonoff universal prior does not eliminate stipulation; it replaces a domain-specific stipulation with a single global stipulation whose effect is bounded on every representable hypothesis simultaneously. Whether that substitution counts as `principled' depends on how much one values universality-up-to-constant over domain fit.'' The Abstract and Introduction now describe the prior as ``universal up to an additive invariance constant.''

\paragraph{MC4 (no consistency/convergence discussion).} Addressed. Section~8.2 states a consistency conjecture for the interactive realizable case, cites Blackwell-Dubins merging and Leike-Hutter convergence as the relevant antecedents, and specifies what an interactive extension would require. Section~8.3 states the exploration hypothesis that makes the expected-free-energy policy sufficient for consistency. These are stated as conjecture and open technical steps rather than as proved theorems, which is the honest position.

\paragraph{MC5 (``open mechanistic challenge'' is aspirational).} Addressed. Section~10 has been substantially expanded. It now specifies, for each of three required internal quantities ($q_t$, $\Gcal_t$, the program neighborhood of $\hat p$), the required access mode, the verification test, and the tolerance. Section~10.2 sketches a statistical-power analysis. Section~10.3 connects the requirements to linear probes, mechanistic interpretability, sparse autoencoders, and ROME-style causal interventions. Section~10.4 articulates compositional compliance and distinguishes the challenge from generic interpretability. Table~3 summarizes the compliance requirements in a form directly usable as an evaluation checklist.

\section*{Response to Minor Technical Concerns}

\begin{description}
\item[T1 (uniqueness in Thm~5.1 and Prop~6.2).] Fixed. Both theorems now include the phrase ``by strict convexity of KL on its effective domain'' in their statements or proofs.
\item[T2 (admissibility and zero-support histories).] Fixed. Section~5 now includes a paragraph immediately after the admissibility definition explaining what happens when $\xi(o_{1:t}\givena a_{1:t}) = 0$ (admissible set empty, variational problem vacuous), and noting that the realizable case avoids this.
\item[T3 (semimeasure-to-distribution transition).] Fixed. Section~7.1 is now devoted to Lemma~7.1 (Semimeasure Completion), which formalizes the $\bot$ absorption and states that all results of Sections~5--6 transfer. The rest of Section~7 works with completed kernels.
\item[T4 (decomposition conditions flagging in Section~6).] Fixed. The proof of Theorem~6.2 (Point-Mass Reduction) now explicitly states that point masses satisfy the stronger integrability conditions of Proposition~5.3, making the decomposition applicable.
\item[T5 (MAP ties and existence).] Fixed. Section~6 now includes a one-line argument for existence (level sets of $J_t^\lambda$ are finite for $\lambda > 0$) and notes that $\argmin$ should be interpreted as a set when ties occur.
\item[T6 (semantic MDL well-definedness).] Fixed. Proposition~6.5 now opens with the statement that $L_t$ depends on $p$ only through $\mu_p$, making the objective well-defined on semantic equivalence classes.
\item[T7 ($\rho_t$ motivation and AIF alignment).] Fixed. Section~7.4 now opens with an explicit reference to \citet{ParrFriston2019}, identifies $\rho_t$ as the goal/target distribution in active inference, and describes two specializations (pure information-seeking and preferred-outcome).
\item[T8 ($\mathcal{P}$ effective enumerability).] Fixed. Section~4.2 now states that $\mathcal{P}$ is effectively enumerable.
\item[T9 (bibliography hygiene).] Fixed. The TODO comment block at the top of the source has been removed; the in-text ``NOTE TO AUTHOR'' has been replaced with finalized citations; the \texttt{.bib} file has been produced; the author list is intentionally blank for anonymized review (with an explicit comment to that effect).
\item[T10 (``exact'' overuse).] Fixed. ``Three exact results'' has been replaced by ``three analytic results'' in the Abstract, Introduction, and Conclusion. ``Missing synthesis'' and ``mathematically legible'' have been replaced with more restrained language (see the rewritten Introduction and Conclusion).
\end{description}

\section*{Response to Reviewer Questions}

\begin{description}
\item[Q1: Convergence or consistency result for $q_t^\star$.] Section~8.2 states a consistency conjecture for the interactive realizable case under an exploration hypothesis, with pointers to Blackwell-Dubins and Leike-Hutter as the relevant antecedents. We flag this as open rather than claim it; the variational-level ``specification'' at each finite history is exact, the asymptotic concentration is conjectural.
\item[Q2: Reference-machine dependence on finite samples.] Section~8.1 now treats this explicitly. Differences within the invariance constant $c_{U,V}$ are not supported by the theory as meaningful ordinal comparisons; see also the Answers-to-Anticipated-Questions subsection in the Discussion.
\item[Q3: Proxy demonstration.] Deferred to a companion paper, as noted under MC2. The revised manuscript fixes all design parameters so that empirical results can be reported against a stable theoretical target.
\item[Q4: $\rho_t$ vs.\ standard AIF target distribution.] Section~7.4 now treats $\rho_t$ explicitly as the goal/target distribution of \citet{ParrFriston2019}. The Answers-to-Anticipated-Questions subsection states: ``Identically in form, once $\rho_t$ is read as the goal/target distribution over next observations.''
\item[Q5: Mechanistic challenge vs.\ general interpretability.] Section~10.4 explicitly distinguishes the two. Generic interpretability asks for human-understandable explanations of model behavior; our mechanistic challenge specifies which internal quantities matter, what verification tests apply to each, and what compositional compliance looks like. Section~10.3 connects the requirements to probing, circuits, SAEs, and causal interventions to show the requirements are not intractable in principle.
\end{description}

\section*{Detailed By-Section Responses}

\begin{description}
\item[Abstract.] Rewritten per MC1 and T10. The phrase ``the missing synthesis'' has been removed and replaced with a more restrained statement.
\item[Introduction.] The ``sharp move'' paragraph has been rewritten. A new paragraph on information-theoretic bounded rationality, compression-based curiosity, and Bayesian experimental design has been added before the contribution statement. The contribution is now described as ``modest but specific.''
\item[Section~2 (Motivation and Explanatory Criteria).] The hard-to-vary criterion now cites \citet{Deutsch2011} as its source.
\item[Section~3 (Related Work).] Restructured into four clusters (algorithmic information theory; FEP; information-theoretic bounded rationality and experimental design; program induction). Each cluster ends with a \emph{Delta} paragraph. A new Positioning subsection (3.5) closes the section.
\item[Section~4 (Mathematical Substrate).] Added one sentence stating effective enumerability of $\mathcal{P}$.
\item[Sections~5--6 (Algorithmic Free Energy, Variational Identity, MDL bridge).] Technical tightening per T1-T6 and admissibility discussion per T2.
\item[Section~7 (Expected Algorithmic Free Energy).] New Section~7.1 contains the Semimeasure Completion Lemma; Section~7.2 disambiguates joint and marginal predictive-law notation; Section~7.4 motivates $\rho_t$ explicitly. Corollary~7.7 now explicitly names Bayesian experimental design and Schmidhuber-style compression curiosity as alternative readings of the deterministic-discovery specialization.
\item[Section~8 (Reference-Machine Invariance, Posterior Consistency, and Exploration).] New section addressing MC3 and MC4.
\item[Section~9 (Behavioral Proxy).] Table~2 is preserved; language around ``honest limitation'' tempered per MC3; an explicit note defers empirical assessment to future work.
\item[Section~10 (The Open Mechanistic Challenge).] Substantially expanded per MC5. Now a compliance protocol with three verification tests, a power analysis, a connection to interpretability tooling, and a requirements table.
\item[Section~11 (Discussion).] Restructured. A new ``Answers to Anticipated Questions'' subsection addresses each of the reviewer's five questions in-line.
\item[Section~12 (Conclusion).] Rewritten to replace ``mathematically legible'' with specific claims and to align the summary with the tempered novelty framing.
\end{description}

\section*{Files in the Revision Package}

\begin{itemize}
\item \texttt{algorithmic\_free\_energy\_principle.tex} --- revised manuscript.
\item \texttt{algorithmic\_free\_energy\_principle.bib} --- updated bibliography.
\item \texttt{response\_letter.tex} (this file) --- point-by-point response.
\end{itemize}

We are grateful to the reviewer for the time and care evident in the report. We believe the revised manuscript is a substantially stronger paper, and we would welcome any further feedback on the changes we have made.

\end{document}
